\documentclass[11pt,a4paper]{article}
\usepackage[margin=2.5cm]{geometry}
\usepackage{amsmath,amssymb}
\usepackage[colorlinks=true,linkcolor=blue]{hyperref}

\title{Analytical Solution of Fission Product Diffusion\\
in a TRISO Particle\\[0.4em]
\large Expanded-Notation Working}
\author{Ray}
\date{\today}

\begin{document}
\maketitle

\begin{abstract}
This document reproduces the Sturm--Liouville analytical solution for transient
fission-product diffusion in a spherical TRISO particle with a constant volumetric
source, zero initial inventory, and a convective (Robin) surface condition, written
throughout in fully expanded notation. Every derivative is written in Leibniz form
(no prime notation), every intermediate algebraic step is retained, and no
dimensionless groupings are introduced: the physical constants $D$, $h$, $R$ and
$S_{0}$ appear explicitly in every expression. The shutdown scenario (source
switched off at $t_{c}$) is included, and the extension to the full five-layer
particle is formulated in the final section.
\end{abstract}

\tableofcontents

\section{Notation}

Two derivative symbols appear in this document, and the distinction between them is
carried consistently.

\begin{center}
\begin{tabular}{llll}
\hline
Function & Depends on & Symbol & Example \\
\hline
$c(r,t)$   & $r$ and $t$ & $\partial$ & $\dfrac{\partial c}{\partial t}$,\quad $\dfrac{\partial c}{\partial r}$ \\[0.9em]
$v(r,t)$   & $r$ and $t$ & $\partial$ & $\dfrac{\partial v}{\partial t}$,\quad $\dfrac{\partial v}{\partial r}$ \\[0.9em]
$w(r)$     & $r$ only    & $d$        & $\dfrac{dw}{dr}$ \\[0.9em]
$\phi(r)$  & $r$ only    & $d$        & $\dfrac{d\phi}{dr}$,\quad $\dfrac{d^{2}\phi}{dr^{2}}$ \\[0.9em]
$u(r)$     & $r$ only    & $d$        & $\dfrac{du}{dr}$,\quad $\dfrac{d^{2}u}{dr^{2}}$ \\[0.9em]
$T(t)$     & $t$ only    & $d$        & $\dfrac{dT}{dt}$ \\
\hline
\end{tabular}
\end{center}

\medskip
\noindent
The symbol $\partial$ signals that other variables are being held fixed. For a
function of a single variable there is nothing to hold fixed, so the ordinary
symbol $d$ is used. The change from $\partial$ to $d$ occurs at
Section~\ref{sec:separation}, and is precisely what separation of variables
achieves: one partial differential equation is replaced by two ordinary
differential equations.

\medskip
\noindent
The notation
\begin{equation*}
\frac{d\phi}{dr}\bigg|_{r=R}
\end{equation*}
means: perform the differentiation first, then substitute $r=R$ into the result.

\section{Problem statement}

Let $c(r,t)$ denote the fission-product concentration
($\mathrm{mol\ m^{-3}}$) in a sphere of radius $R$ with diffusivity $D$
($\mathrm{m^{2}\ s^{-1}}$), constant volumetric source $S_{0}$
($\mathrm{mol\ m^{-3}\ s^{-1}}$), and surface mass-transfer coefficient $h$
($\mathrm{m\ s^{-1}}$).

\begin{equation}
\frac{\partial c}{\partial t}
= \frac{D}{r^{2}}\,\frac{\partial}{\partial r}\!\left(r^{2}\frac{\partial c}{\partial r}\right) + S_{0},
\qquad 0 < r < R,\quad t > 0 .
\label{eq:pde}
\end{equation}

Expanding the inner derivative by the product rule gives the equivalent form
\begin{equation}
\frac{\partial c}{\partial t}
= D\left(\frac{\partial^{2} c}{\partial r^{2}} + \frac{2}{r}\,\frac{\partial c}{\partial r}\right) + S_{0} .
\label{eq:pde-expanded}
\end{equation}

The initial and boundary data are
\begin{align}
c(r,0) &= 0 , \label{eq:ic}\\[0.3em]
\frac{\partial c}{\partial r}\bigg|_{r=0} &= 0
\qquad \text{(spherical symmetry)} , \label{eq:bc0}\\[0.3em]
-D\,\frac{\partial c}{\partial r}\bigg|_{r=R} &= h\,c(R,t)
\qquad \text{(convective release)} . \label{eq:bcR}
\end{align}

The source is active for $0 \le t \le t_{c}$ and is switched off for $t > t_{c}$.

\section{Steady--transient split}

\begin{equation}
c(r,t) = w(r) + v(r,t) .
\label{eq:split}
\end{equation}

\section{Steady solution}

Setting the time derivative to zero in \eqref{eq:pde} leaves a function of $r$
alone, so ordinary derivatives apply from here on.

\begin{equation}
\frac{D}{r^{2}}\,\frac{d}{dr}\!\left(r^{2}\frac{dw}{dr}\right) + S_{0} = 0 .
\end{equation}

Multiplying through by $r^{2}/D$,
\begin{equation}
\frac{d}{dr}\!\left(r^{2}\frac{dw}{dr}\right) = -\frac{S_{0}\,r^{2}}{D} .
\end{equation}

Integrating once with respect to $r$, with $A$ the constant of integration,
\begin{equation}
r^{2}\frac{dw}{dr} = -\frac{S_{0}\,r^{3}}{3D} + A .
\end{equation}

Dividing through by $r^{2}$,
\begin{equation}
\frac{dw}{dr} = -\frac{S_{0}\,r}{3D} + \frac{A}{r^{2}} .
\end{equation}

The term $A/r^{2}$ diverges as $r \to 0$, which would give unbounded flux at the
centre and contradict \eqref{eq:bc0}. Hence
\begin{equation}
A = 0 ,
\qquad\text{so}\qquad
\frac{dw}{dr} = -\frac{S_{0}\,r}{3D} .
\end{equation}

Integrating a second time, with $B$ the second constant of integration,
\begin{equation}
w(r) = -\frac{S_{0}\,r^{2}}{6D} + B .
\end{equation}

Evaluating the gradient at the surface,
\begin{equation}
\frac{dw}{dr}\bigg|_{r=R} = -\frac{S_{0}R}{3D} ,
\end{equation}
and substituting into the surface condition \eqref{eq:bcR},
\begin{equation}
-D\,\frac{dw}{dr}\bigg|_{r=R}
= -D\left(-\frac{S_{0}R}{3D}\right)
= \frac{S_{0}R}{3}
= h\,w(R)
= h\left(B - \frac{S_{0}R^{2}}{6D}\right) .
\end{equation}

Solving for $B$,
\begin{equation}
B = \frac{S_{0}R}{3h} + \frac{S_{0}R^{2}}{6D} ,
\end{equation}
so that
\begin{equation}
\boxed{\;
w(r) = \frac{S_{0}R}{3h} + \frac{S_{0}}{6D}\left(R^{2} - r^{2}\right) .
\;}
\label{eq:w}
\end{equation}

\section{Transient problem}

Subtracting the steady part removes the source and homogenises both boundary
conditions:
\begin{equation}
\frac{\partial v}{\partial t}
= \frac{D}{r^{2}}\,\frac{\partial}{\partial r}\!\left(r^{2}\frac{\partial v}{\partial r}\right) ,
\label{eq:vpde}
\end{equation}
\begin{equation}
\frac{\partial v}{\partial r}\bigg|_{r=0} = 0 ,
\qquad
-D\,\frac{\partial v}{\partial r}\bigg|_{r=R} = h\,v(R,t) ,
\qquad
v(r,0) = -w(r) .
\label{eq:vdata}
\end{equation}

\section{Separation of variables}
\label{sec:separation}

Seek a solution in the form of a pure product,
\begin{equation}
v(r,t) = \phi(r)\,T(t) .
\end{equation}

Differentiating with respect to $t$, with $r$ held fixed, the factor $\phi(r)$
contains no $t$ and behaves as a constant multiplier:
\begin{equation}
\frac{\partial v}{\partial t} = \phi(r)\,\frac{dT}{dt} .
\end{equation}

Differentiating with respect to $r$, with $t$ held fixed, the factor $T(t)$
contains no $r$ and behaves as a constant multiplier:
\begin{equation}
\frac{\partial v}{\partial r} = \frac{d\phi}{dr}\,T(t) .
\end{equation}

Multiplying by $r^{2}$,
\begin{equation}
r^{2}\frac{\partial v}{\partial r} = r^{2}\frac{d\phi}{dr}\,T(t) .
\end{equation}

Differentiating once more with respect to $r$, again factoring $T(t)$ out of the
$r$-derivative,
\begin{equation}
\frac{\partial}{\partial r}\!\left(r^{2}\frac{\partial v}{\partial r}\right)
= \frac{d}{dr}\!\left(r^{2}\frac{d\phi}{dr}\right) T(t) .
\label{eq:sep-key}
\end{equation}

Equation \eqref{eq:sep-key} is where the notation changes: the quantity
$r^{2}\,d\phi/dr$ contains no $t$ whatsoever, so the derivative acting on it is
ordinary rather than partial.

Substituting into \eqref{eq:vpde},
\begin{equation}
\phi(r)\,\frac{dT}{dt}
= \frac{D}{r^{2}}\,\frac{d}{dr}\!\left(r^{2}\frac{d\phi}{dr}\right) T(t) .
\end{equation}

Dividing both sides by $D\,\phi(r)\,T(t)$,
\begin{equation}
\frac{1}{D\,T(t)}\,\frac{dT}{dt}
= \frac{1}{r^{2}\phi(r)}\,\frac{d}{dr}\!\left(r^{2}\frac{d\phi}{dr}\right) .
\label{eq:separated}
\end{equation}

The left-hand side of \eqref{eq:separated} depends only on $t$; the right-hand side
depends only on $r$. Since $r$ and $t$ vary independently, both sides must equal the
same constant, written $-\lambda$:
\begin{equation}
\frac{1}{D\,T(t)}\,\frac{dT}{dt}
= \frac{1}{r^{2}\phi(r)}\,\frac{d}{dr}\!\left(r^{2}\frac{d\phi}{dr}\right)
= -\lambda .
\end{equation}

This yields two ordinary differential equations. The time equation,
\begin{equation}
\frac{dT}{dt} = -\lambda D\,T(t)
\qquad \Longrightarrow \qquad
T(t) = e^{-\lambda D t} ,
\label{eq:timeeq}
\end{equation}
and the spatial equation,
\begin{equation}
\frac{d}{dr}\!\left(r^{2}\frac{d\phi}{dr}\right) + \lambda\,r^{2}\phi = 0 .
\label{eq:spatial}
\end{equation}

\section{Removal of the curvature term}

Introduce
\begin{equation}
u(r) = r\,\phi(r) ,
\qquad\text{equivalently}\qquad
\phi(r) = \frac{u(r)}{r} .
\end{equation}

Differentiating $\phi = u \cdot r^{-1}$ by the product rule,
\begin{equation}
\frac{d\phi}{dr}
= \frac{1}{r}\,\frac{du}{dr} + u \cdot \left(-\frac{1}{r^{2}}\right)
= \frac{1}{r}\,\frac{du}{dr} - \frac{u}{r^{2}} .
\end{equation}

Multiplying through by $r^{2}$,
\begin{equation}
r^{2}\frac{d\phi}{dr} = r\,\frac{du}{dr} - u .
\label{eq:r2phi}
\end{equation}

Differentiating \eqref{eq:r2phi} with respect to $r$, applying the product rule to
the first term,
\begin{equation}
\frac{d}{dr}\!\left(r^{2}\frac{d\phi}{dr}\right)
= \frac{d}{dr}\!\left(r\,\frac{du}{dr} - u\right)
= \left(\frac{du}{dr} + r\,\frac{d^{2}u}{dr^{2}}\right) - \frac{du}{dr}
= r\,\frac{d^{2}u}{dr^{2}} .
\label{eq:key-identity}
\end{equation}

The remaining term of \eqref{eq:spatial} becomes
\begin{equation}
r^{2}\phi = r^{2}\cdot\frac{u}{r} = r\,u .
\end{equation}

Substituting both results into \eqref{eq:spatial},
\begin{equation}
r\,\frac{d^{2}u}{dr^{2}} + \lambda\,r\,u = 0 ,
\end{equation}
and dividing by $r$ (valid for $r > 0$),
\begin{equation}
\boxed{\;
\frac{d^{2}u}{dr^{2}} + \lambda u = 0 .
\;}
\label{eq:oscillator}
\end{equation}

All curvature terms have cancelled.

\section{General solution and regularity at the origin}

Substituting the trial form $u(r) = e^{mr}$ into \eqref{eq:oscillator}, for which
$d^{2}u/dr^{2} = m^{2}e^{mr}$,
\begin{equation}
m^{2} e^{mr} + \lambda e^{mr} = 0
\qquad \Longrightarrow \qquad
m^{2} + \lambda = 0
\qquad \Longrightarrow \qquad
m = \pm\, i\sqrt{\lambda} .
\end{equation}

The corresponding real general solution is
\begin{equation}
u(r) = C_{1}\sin\!\left(\sqrt{\lambda}\,r\right) + C_{2}\cos\!\left(\sqrt{\lambda}\,r\right) ,
\end{equation}
so that
\begin{equation}
\phi(r)
= C_{1}\,\frac{\sin\!\left(\sqrt{\lambda}\,r\right)}{r}
+ C_{2}\,\frac{\cos\!\left(\sqrt{\lambda}\,r\right)}{r} .
\end{equation}

Expanding each branch near the origin,
\begin{equation}
\frac{\sin\!\left(\sqrt{\lambda}\,r\right)}{r}
= \sqrt{\lambda} - \frac{\lambda^{3/2} r^{2}}{6} + \cdots
\;\xrightarrow[\;r \to 0\;]{}\;
\sqrt{\lambda}
\qquad \text{(finite)} ,
\end{equation}
\begin{equation}
\frac{\cos\!\left(\sqrt{\lambda}\,r\right)}{r}
= \frac{1}{r} - \frac{\lambda r}{2} + \cdots
\;\xrightarrow[\;r \to 0\;]{}\;
\infty
\qquad \text{(unbounded)} .
\end{equation}

Boundedness of the concentration at the centre therefore requires
\begin{equation}
C_{2} = 0 ,
\end{equation}
and setting $C_{1} = 1$ (any overall constant is absorbed later into the expansion
coefficients),
\begin{equation}
\boxed{\;
u(r) = \sin\!\left(\sqrt{\lambda}\,r\right) ,
\qquad
\phi(r) = \frac{\sin\!\left(\sqrt{\lambda}\,r\right)}{r} .
\;}
\label{eq:eigfun}
\end{equation}

The centre condition \eqref{eq:bc0} is then satisfied automatically, since
\begin{equation}
\phi(r) = \sqrt{\lambda} - \frac{\lambda^{3/2} r^{2}}{6} + \cdots
\qquad \Longrightarrow \qquad
\frac{d\phi}{dr} = -\frac{\lambda^{3/2} r}{3} + \cdots
\qquad \Longrightarrow \qquad
\frac{d\phi}{dr}\bigg|_{r=0} = 0 .
\end{equation}

\section{Eigencondition}

Evaluating \eqref{eq:eigfun} at the surface,
\begin{equation}
\phi(R) = \frac{\sin\!\left(\sqrt{\lambda}\,R\right)}{R} .
\end{equation}

Differentiating $\phi(r) = \sin\!\left(\sqrt{\lambda}\,r\right) \cdot r^{-1}$ by the
product rule,
\begin{equation}
\frac{d\phi}{dr}
= \sqrt{\lambda}\cos\!\left(\sqrt{\lambda}\,r\right)\cdot\frac{1}{r}
- \sin\!\left(\sqrt{\lambda}\,r\right)\cdot\frac{1}{r^{2}} ,
\end{equation}
and combining over a common denominator,
\begin{equation}
\frac{d\phi}{dr}
= \frac{\sqrt{\lambda}\,r\cos\!\left(\sqrt{\lambda}\,r\right)
- \sin\!\left(\sqrt{\lambda}\,r\right)}{r^{2}} .
\end{equation}

At the surface,
\begin{equation}
\frac{d\phi}{dr}\bigg|_{r=R}
= \frac{\sqrt{\lambda}\,R\cos\!\left(\sqrt{\lambda}\,R\right)
- \sin\!\left(\sqrt{\lambda}\,R\right)}{R^{2}} .
\end{equation}

Substituting into the Robin condition
$-D\,(d\phi/dr)|_{r=R} = h\,\phi(R)$,
\begin{equation}
-D \cdot \frac{\sqrt{\lambda}\,R\cos\!\left(\sqrt{\lambda}\,R\right)
- \sin\!\left(\sqrt{\lambda}\,R\right)}{R^{2}}
= h \cdot \frac{\sin\!\left(\sqrt{\lambda}\,R\right)}{R} .
\end{equation}

Multiplying both sides by $R^{2}$,
\begin{equation}
-D\left(\sqrt{\lambda}\,R\cos\!\left(\sqrt{\lambda}\,R\right)
- \sin\!\left(\sqrt{\lambda}\,R\right)\right)
= hR\sin\!\left(\sqrt{\lambda}\,R\right) .
\end{equation}

Distributing the minus sign,
\begin{equation}
D\sin\!\left(\sqrt{\lambda}\,R\right)
- D\sqrt{\lambda}\,R\cos\!\left(\sqrt{\lambda}\,R\right)
= hR\sin\!\left(\sqrt{\lambda}\,R\right) ,
\end{equation}
and collecting the sine terms,
\begin{equation}
\left(D - hR\right)\sin\!\left(\sqrt{\lambda}\,R\right)
= D\,\sqrt{\lambda}\,R\cos\!\left(\sqrt{\lambda}\,R\right) .
\label{eq:eig-precursor}
\end{equation}

Dividing by $\sin\!\left(\sqrt{\lambda}\,R\right)$ gives the eigencondition,
\begin{equation}
\boxed{\;
D\,\sqrt{\lambda}\,R\cot\!\left(\sqrt{\lambda}\,R\right) = D - hR .
\;}
\label{eq:eigcond}
\end{equation}

Equation \eqref{eq:eigcond} is transcendental and admits infinitely many roots,
\begin{equation}
\lambda_{1} < \lambda_{2} < \lambda_{3} < \cdots ,
\qquad
(n-1)\pi < \sqrt{\lambda_{n}}\,R < n\pi ,
\end{equation}
with corresponding eigenfunctions
\begin{equation}
\phi_{n}(r) = \frac{\sin\!\left(\sqrt{\lambda_{n}}\,r\right)}{r} .
\end{equation}

Rearranging \eqref{eq:eig-precursor} gives the identity used repeatedly below to
eliminate the cosine:
\begin{equation}
\cos\!\left(\sqrt{\lambda_{n}}\,R\right)
= \frac{\left(D - hR\right)\sin\!\left(\sqrt{\lambda_{n}}\,R\right)}
{D\,\sqrt{\lambda_{n}}\,R} .
\label{eq:coselim}
\end{equation}

\section{Positivity of the eigenvalues}

Multiplying \eqref{eq:spatial} by $\phi$ and integrating over the sphere,
\begin{equation}
\int_{0}^{R} \phi\,\frac{d}{dr}\!\left(r^{2}\frac{d\phi}{dr}\right) dr
+ \lambda \int_{0}^{R} r^{2}\phi^{2}\,dr = 0 .
\end{equation}

The first integral is treated by parts. Since
\begin{equation}
\frac{d}{dr}\!\left(\phi \cdot r^{2}\frac{d\phi}{dr}\right)
= \frac{d\phi}{dr}\cdot r^{2}\frac{d\phi}{dr}
+ \phi\,\frac{d}{dr}\!\left(r^{2}\frac{d\phi}{dr}\right) ,
\end{equation}
it follows that
\begin{equation}
\int_{0}^{R} \phi\,\frac{d}{dr}\!\left(r^{2}\frac{d\phi}{dr}\right) dr
= \left[r^{2}\phi\,\frac{d\phi}{dr}\right]_{0}^{R}
- \int_{0}^{R} r^{2}\left(\frac{d\phi}{dr}\right)^{2} dr .
\end{equation}

Substituting and solving for $\lambda$,
\begin{equation}
\lambda =
\frac{\displaystyle -\left[r^{2}\phi\,\frac{d\phi}{dr}\right]_{0}^{R}
+ \int_{0}^{R} r^{2}\left(\frac{d\phi}{dr}\right)^{2} dr}
{\displaystyle \int_{0}^{R} r^{2}\phi^{2}\,dr} .
\label{eq:rayleigh}
\end{equation}

At the lower limit the factor $r^{2}$ vanishes,
\begin{equation}
\left[r^{2}\phi\,\frac{d\phi}{dr}\right]_{r=0} = 0 ,
\end{equation}
while at the upper limit the Robin condition gives
\begin{equation}
\frac{d\phi}{dr}\bigg|_{r=R} = -\frac{h}{D}\,\phi(R) ,
\end{equation}
so that
\begin{equation}
-\left[r^{2}\phi\,\frac{d\phi}{dr}\right]_{r=R}
= -R^{2}\phi(R)\left(-\frac{h}{D}\phi(R)\right)
= \frac{h}{D}\,R^{2}\phi(R)^{2} \;\ge\; 0 .
\end{equation}

Both terms in the numerator of \eqref{eq:rayleigh} are non-negative and the
denominator is strictly positive, so $\lambda \ge 0$. The value $\lambda = 0$ would
force $r^{2}\,d\phi/dr$ constant, hence zero by regularity, hence $\phi$ constant;
the Robin condition with $h > 0$ then forces $\phi \equiv 0$, which is not an
eigenfunction. Therefore
\begin{equation}
\lambda_{n} > 0 \qquad \text{for all } n ,
\end{equation}
and every mode $e^{-\lambda_{n} D t}$ decays.

\section{Orthogonality}

Let $\phi_{n}$ and $\phi_{m}$ be eigenfunctions with $\lambda_{n} \neq \lambda_{m}$:
\begin{align}
\frac{d}{dr}\!\left(r^{2}\frac{d\phi_{n}}{dr}\right) + \lambda_{n} r^{2}\phi_{n} &= 0 , \\
\frac{d}{dr}\!\left(r^{2}\frac{d\phi_{m}}{dr}\right) + \lambda_{m} r^{2}\phi_{m} &= 0 .
\end{align}

Multiplying the first by $\phi_{m}$, the second by $\phi_{n}$, and subtracting,
\begin{equation}
\phi_{m}\,\frac{d}{dr}\!\left(r^{2}\frac{d\phi_{n}}{dr}\right)
- \phi_{n}\,\frac{d}{dr}\!\left(r^{2}\frac{d\phi_{m}}{dr}\right)
= \left(\lambda_{m} - \lambda_{n}\right) r^{2}\phi_{n}\phi_{m} .
\end{equation}

The left-hand side is an exact derivative, since the cross terms cancel:
\begin{equation}
\frac{d}{dr}\!\left[r^{2}\!\left(\phi_{m}\frac{d\phi_{n}}{dr}
- \phi_{n}\frac{d\phi_{m}}{dr}\right)\right]
= \phi_{m}\,\frac{d}{dr}\!\left(r^{2}\frac{d\phi_{n}}{dr}\right)
- \phi_{n}\,\frac{d}{dr}\!\left(r^{2}\frac{d\phi_{m}}{dr}\right) .
\end{equation}

Integrating from $0$ to $R$, the boundary term vanishes at the origin because of the
factor $r^{2}$,
\begin{equation}
\left[r^{2}\!\left(\phi_{m}\frac{d\phi_{n}}{dr}
- \phi_{n}\frac{d\phi_{m}}{dr}\right)\right]_{r=0} = 0 ,
\end{equation}
and vanishes at the surface because both eigenfunctions obey the same Robin
relation,
\begin{equation}
\left[\phi_{m}\frac{d\phi_{n}}{dr} - \phi_{n}\frac{d\phi_{m}}{dr}\right]_{r=R}
= \phi_{m}\left(-\frac{h}{D}\phi_{n}\right)
- \phi_{n}\left(-\frac{h}{D}\phi_{m}\right) = 0 .
\end{equation}

Hence
\begin{equation}
\left(\lambda_{m} - \lambda_{n}\right)\int_{0}^{R} r^{2}\phi_{n}\phi_{m}\,dr = 0 ,
\end{equation}
and since $\lambda_{m} \neq \lambda_{n}$,
\begin{equation}
\boxed{\;
\int_{0}^{R} \phi_{n}\,\phi_{m}\, r^{2}\,dr = 0 ,
\qquad n \neq m .
\;}
\label{eq:orthog}
\end{equation}

\section{Normalization}

With $\phi_{n}^{2} r^{2} = \sin^{2}\!\left(\sqrt{\lambda_{n}}\,r\right)$,
\begin{equation}
N_{n} = \int_{0}^{R} \phi_{n}^{2}\, r^{2}\,dr
= \int_{0}^{R} \sin^{2}\!\left(\sqrt{\lambda_{n}}\,r\right) dr .
\end{equation}

Applying the half-angle identity,
\begin{equation}
N_{n}
= \int_{0}^{R} \frac{1 - \cos\!\left(2\sqrt{\lambda_{n}}\,r\right)}{2}\,dr
= \frac{R}{2} - \frac{\sin\!\left(2\sqrt{\lambda_{n}}\,R\right)}{4\sqrt{\lambda_{n}}} .
\end{equation}

Using $\sin(2\theta) = 2\sin\theta\cos\theta$,
\begin{equation}
N_{n} = \frac{R}{2}
- \frac{\sin\!\left(\sqrt{\lambda_{n}}\,R\right)
\cos\!\left(\sqrt{\lambda_{n}}\,R\right)}{2\sqrt{\lambda_{n}}} ,
\end{equation}
and eliminating the cosine with \eqref{eq:coselim},
\begin{equation}
N_{n} = \frac{R}{2}
- \frac{\left(D - hR\right)\sin^{2}\!\left(\sqrt{\lambda_{n}}\,R\right)}
{2D\,\lambda_{n} R} ,
\end{equation}
so that
\begin{equation}
\boxed{\;
N_{n} = \frac{R}{2}\left[1
- \frac{\left(D - hR\right)\sin^{2}\!\left(\sqrt{\lambda_{n}}\,R\right)}
{D\,\lambda_{n} R^{2}}\right] .
\;}
\label{eq:norm}
\end{equation}

\section{Projection of the initial condition}

Sturm--Liouville theory guarantees that the eigenfunctions form a complete basis, so
\begin{equation}
v(r,0) = \sum_{n \ge 1} a_{n}\,\phi_{n}(r) .
\end{equation}

Multiplying by $\phi_{n} r^{2}$ and integrating, all cross terms vanish by
\eqref{eq:orthog}, leaving
\begin{equation}
\int_{0}^{R} v(r,0)\,\phi_{n}(r)\, r^{2}\,dr = a_{n} N_{n} ,
\qquad\text{hence}\qquad
a_{n} = \frac{1}{N_{n}}\int_{0}^{R} v(r,0)\,\phi_{n}(r)\, r^{2}\,dr .
\end{equation}

Using $v(r,0) = -w(r)$ from \eqref{eq:vdata} and
$\phi_{n}(r)\,r^{2} = r\sin\!\left(\sqrt{\lambda_{n}}\,r\right)$,
\begin{equation}
a_{n} = -\frac{1}{N_{n}}
\int_{0}^{R} w(r)\, r\sin\!\left(\sqrt{\lambda_{n}}\,r\right) dr .
\label{eq:an-def}
\end{equation}

Writing $w(r)$ from \eqref{eq:w} in the form of a constant minus a quadratic,
\begin{equation}
w(r) = \left(\frac{S_{0}R}{3h} + \frac{S_{0}R^{2}}{6D}\right)
- \frac{S_{0}}{6D}\,r^{2} ,
\end{equation}
the required integral separates as
\begin{equation}
\int_{0}^{R} w(r)\, r\sin\!\left(\sqrt{\lambda_{n}}\,r\right) dr
= \left(\frac{S_{0}R}{3h} + \frac{S_{0}R^{2}}{6D}\right) I_{1}
- \frac{S_{0}}{6D}\, I_{3} ,
\label{eq:wsplit}
\end{equation}
where
\begin{equation}
I_{1} = \int_{0}^{R} r\sin\!\left(\sqrt{\lambda_{n}}\,r\right) dr ,
\qquad
I_{3} = \int_{0}^{R} r^{3}\sin\!\left(\sqrt{\lambda_{n}}\,r\right) dr .
\end{equation}

\section{Evaluation of the integrals}

\subsection{The integral \texorpdfstring{$I_{1}$}{I1}}

Integrating by parts with $u = r$ and
$dv = \sin\!\left(\sqrt{\lambda_{n}}\,r\right) dr$,
\begin{equation}
I_{1}
= \left[-\frac{r\cos\!\left(\sqrt{\lambda_{n}}\,r\right)}{\sqrt{\lambda_{n}}}\right]_{0}^{R}
+ \frac{1}{\sqrt{\lambda_{n}}}\int_{0}^{R}\cos\!\left(\sqrt{\lambda_{n}}\,r\right) dr .
\end{equation}

Evaluating both pieces,
\begin{equation}
I_{1}
= -\frac{R\cos\!\left(\sqrt{\lambda_{n}}\,R\right)}{\sqrt{\lambda_{n}}}
+ \frac{1}{\sqrt{\lambda_{n}}}
\left[\frac{\sin\!\left(\sqrt{\lambda_{n}}\,r\right)}{\sqrt{\lambda_{n}}}\right]_{0}^{R}
= \frac{\sin\!\left(\sqrt{\lambda_{n}}\,R\right)
- \sqrt{\lambda_{n}}\,R\cos\!\left(\sqrt{\lambda_{n}}\,R\right)}{\lambda_{n}} .
\end{equation}

From \eqref{eq:coselim},
\begin{equation}
\sqrt{\lambda_{n}}\,R\cos\!\left(\sqrt{\lambda_{n}}\,R\right)
= \frac{\left(D - hR\right)\sin\!\left(\sqrt{\lambda_{n}}\,R\right)}{D} ,
\end{equation}
so that
\begin{equation}
I_{1}
= \frac{\sin\!\left(\sqrt{\lambda_{n}}\,R\right)}{\lambda_{n}}
\left[1 - \frac{D - hR}{D}\right] ,
\end{equation}
giving
\begin{equation}
\boxed{\;
I_{1} = \frac{hR\,\sin\!\left(\sqrt{\lambda_{n}}\,R\right)}{D\,\lambda_{n}} .
\;}
\label{eq:I1}
\end{equation}

\subsection{The integral \texorpdfstring{$I_{3}$}{I3}}

Integrating by parts with $u = r^{3}$,
\begin{equation}
\int r^{3}\sin\!\left(\sqrt{\lambda_{n}}\,r\right) dr
= -\frac{r^{3}\cos\!\left(\sqrt{\lambda_{n}}\,r\right)}{\sqrt{\lambda_{n}}}
+ \frac{3}{\sqrt{\lambda_{n}}}\int r^{2}\cos\!\left(\sqrt{\lambda_{n}}\,r\right) dr .
\end{equation}

The remaining integral is reduced by a further integration by parts with $u = r^{2}$,
\begin{equation}
\int r^{2}\cos\!\left(\sqrt{\lambda_{n}}\,r\right) dr
= \frac{r^{2}\sin\!\left(\sqrt{\lambda_{n}}\,r\right)}{\sqrt{\lambda_{n}}}
- \frac{2}{\sqrt{\lambda_{n}}}\int r\sin\!\left(\sqrt{\lambda_{n}}\,r\right) dr ,
\end{equation}
and the innermost integral is
\begin{equation}
\int r\sin\!\left(\sqrt{\lambda_{n}}\,r\right) dr
= -\frac{r\cos\!\left(\sqrt{\lambda_{n}}\,r\right)}{\sqrt{\lambda_{n}}}
+ \frac{\sin\!\left(\sqrt{\lambda_{n}}\,r\right)}{\lambda_{n}} .
\end{equation}

Back-substituting,
\begin{equation}
\int r^{2}\cos\!\left(\sqrt{\lambda_{n}}\,r\right) dr
= \frac{r^{2}\sin\!\left(\sqrt{\lambda_{n}}\,r\right)}{\sqrt{\lambda_{n}}}
+ \frac{2r\cos\!\left(\sqrt{\lambda_{n}}\,r\right)}{\lambda_{n}}
- \frac{2\sin\!\left(\sqrt{\lambda_{n}}\,r\right)}{\lambda_{n}^{3/2}} ,
\end{equation}
and therefore
\begin{equation}
\int r^{3}\sin\!\left(\sqrt{\lambda_{n}}\,r\right) dr
= -\frac{r^{3}\cos\!\left(\sqrt{\lambda_{n}}\,r\right)}{\sqrt{\lambda_{n}}}
+ \frac{3r^{2}\sin\!\left(\sqrt{\lambda_{n}}\,r\right)}{\lambda_{n}}
+ \frac{6r\cos\!\left(\sqrt{\lambda_{n}}\,r\right)}{\lambda_{n}^{3/2}}
- \frac{6\sin\!\left(\sqrt{\lambda_{n}}\,r\right)}{\lambda_{n}^{2}} .
\end{equation}

Evaluating between $0$ and $R$, where every term vanishes at the lower limit, and
grouping the sine and cosine contributions,
\begin{equation}
I_{3}
= \left(\frac{3R^{2}}{\lambda_{n}} - \frac{6}{\lambda_{n}^{2}}\right)
\sin\!\left(\sqrt{\lambda_{n}}\,R\right)
- \left(\frac{R^{3}}{\sqrt{\lambda_{n}}} - \frac{6R}{\lambda_{n}^{3/2}}\right)
\cos\!\left(\sqrt{\lambda_{n}}\,R\right) .
\end{equation}

Eliminating the cosine with \eqref{eq:coselim},
\begin{equation}
- \left(\frac{R^{3}}{\sqrt{\lambda_{n}}} - \frac{6R}{\lambda_{n}^{3/2}}\right)
\cos\!\left(\sqrt{\lambda_{n}}\,R\right)
= -\frac{\left(D - hR\right)\sin\!\left(\sqrt{\lambda_{n}}\,R\right)}{D}
\left(\frac{R^{2}}{\lambda_{n}} - \frac{6}{\lambda_{n}^{2}}\right) ,
\end{equation}
so that
\begin{equation}
I_{3} = \sin\!\left(\sqrt{\lambda_{n}}\,R\right)
\left[\frac{R^{2}}{\lambda_{n}}\left(3 - \frac{D - hR}{D}\right)
- \frac{6}{\lambda_{n}^{2}}\left(1 - \frac{D - hR}{D}\right)\right] .
\end{equation}

Simplifying the two bracketed coefficients,
\begin{equation}
3 - \frac{D - hR}{D} = \frac{2D + hR}{D} ,
\qquad
1 - \frac{D - hR}{D} = \frac{hR}{D} ,
\end{equation}
gives
\begin{equation}
\boxed{\;
I_{3} = \sin\!\left(\sqrt{\lambda_{n}}\,R\right)
\left[\frac{\left(2D + hR\right)R^{2}}{D\,\lambda_{n}}
- \frac{6hR}{D\,\lambda_{n}^{2}}\right] .
\;}
\label{eq:I3}
\end{equation}

\section{Expansion coefficients}

Substituting \eqref{eq:I1} and \eqref{eq:I3} into \eqref{eq:wsplit} and grouping by
powers of $\lambda_{n}$,
\begin{align}
\int_{0}^{R} w(r)\, r\sin\!\left(\sqrt{\lambda_{n}}\,r\right) dr
&= \frac{\sin\!\left(\sqrt{\lambda_{n}}\,R\right)}{\lambda_{n}}
\left[\left(\frac{S_{0}R}{3h} + \frac{S_{0}R^{2}}{6D}\right)\frac{hR}{D}
- \frac{S_{0}}{6D}\cdot\frac{\left(2D + hR\right)R^{2}}{D}\right] \notag \\
&\quad + \frac{S_{0}}{6D}\cdot
\frac{6hR\,\sin\!\left(\sqrt{\lambda_{n}}\,R\right)}{D\,\lambda_{n}^{2}} .
\label{eq:grouped}
\end{align}

The bracketed term vanishes identically. Expanding its first part,
\begin{equation}
\left(\frac{S_{0}R}{3h} + \frac{S_{0}R^{2}}{6D}\right)\frac{hR}{D}
= \frac{S_{0}R^{2}}{3D} + \frac{S_{0}hR^{3}}{6D^{2}} ,
\end{equation}
and its second part,
\begin{equation}
\frac{S_{0}}{6D}\cdot\frac{\left(2D + hR\right)R^{2}}{D}
= \frac{2S_{0}R^{2}D}{6D^{2}} + \frac{S_{0}hR^{3}}{6D^{2}}
= \frac{S_{0}R^{2}}{3D} + \frac{S_{0}hR^{3}}{6D^{2}} ,
\end{equation}
so that
\begin{equation}
\left(\frac{S_{0}R}{3h} + \frac{S_{0}R^{2}}{6D}\right)\frac{hR}{D}
- \frac{S_{0}}{6D}\cdot\frac{\left(2D + hR\right)R^{2}}{D} = 0 .
\end{equation}

This cancellation occurs only when the steady state \eqref{eq:w} and the
eigencondition \eqref{eq:eigcond} are mutually consistent. What remains of
\eqref{eq:grouped} is
\begin{equation}
\int_{0}^{R} w(r)\, r\sin\!\left(\sqrt{\lambda_{n}}\,r\right) dr
= \frac{S_{0}}{6D}\cdot
\frac{6hR\,\sin\!\left(\sqrt{\lambda_{n}}\,R\right)}{D\,\lambda_{n}^{2}}
= \frac{hR\,S_{0}\sin\!\left(\sqrt{\lambda_{n}}\,R\right)}{D^{2}\,\lambda_{n}^{2}} ,
\end{equation}
and substituting into \eqref{eq:an-def},
\begin{equation}
\boxed{\;
a_{n} = -\frac{hR\,S_{0}\,\sin\!\left(\sqrt{\lambda_{n}}\,R\right)}
{D^{2}\,\lambda_{n}^{2}\,N_{n}} .
\;}
\label{eq:an}
\end{equation}

\section{Solution with the source active}

Reassembling $c = w + v$ with \eqref{eq:w}, \eqref{eq:eigfun}, \eqref{eq:timeeq} and
\eqref{eq:an},
\begin{equation}
\boxed{\;
c(r,t) = \frac{S_{0}R}{3h} + \frac{S_{0}}{6D}\left(R^{2} - r^{2}\right)
+ \sum_{n \ge 1} a_{n}\,
\frac{\sin\!\left(\sqrt{\lambda_{n}}\,r\right)}{r}\;
e^{-\lambda_{n} D t} ,
\qquad 0 \le t \le t_{c} .
\;}
\label{eq:phase1}
\end{equation}

At $t = 0$ every exponential equals one and the series reproduces $-w(r)$ by
construction, so the initial condition \eqref{eq:ic} is met exactly. As
$t \to \infty$ every mode decays and the profile relaxes to the steady parabola
\eqref{eq:w}.

The instantaneous release rate through the surface is
\begin{equation}
\dot{m}(t) = 4\pi R^{2}\, h\, c(R,t)
\qquad \left(\mathrm{mol\ s^{-1}}\right) ,
\end{equation}
and the slowest mode sets the characteristic equilibration time
\begin{equation}
\tau_{1} = \frac{1}{\lambda_{1} D} .
\end{equation}

\section{Solution after source switch-off}

For $t > t_{c}$ the governing equation is homogeneous with unchanged boundary
conditions,
\begin{equation}
\frac{\partial c}{\partial t}
= \frac{D}{r^{2}}\,\frac{\partial}{\partial r}\!\left(r^{2}\frac{\partial c}{\partial r}\right) ,
\qquad t > t_{c} ,
\end{equation}
so the same eigenbasis applies:
\begin{equation}
c(r,t) = \sum_{n \ge 1} b_{n}\,
\frac{\sin\!\left(\sqrt{\lambda_{n}}\,r\right)}{r}\;
e^{-\lambda_{n} D\,(t - t_{c})} .
\end{equation}

The new initial state is the phase-one solution evaluated at $t_{c}$. Since the
projection of $w$ onto the eigenbasis is $-a_{n}$ by the definition
\eqref{eq:an-def},
\begin{equation}
w(r) = -\sum_{n \ge 1} a_{n}\,\phi_{n}(r) ,
\end{equation}
and therefore
\begin{equation}
c(r,t_{c})
= w(r) + \sum_{n \ge 1} a_{n}\,\phi_{n}(r)\, e^{-\lambda_{n} D t_{c}}
= \sum_{n \ge 1}\left(-a_{n} + a_{n} e^{-\lambda_{n} D t_{c}}\right)\phi_{n}(r) ,
\end{equation}
giving
\begin{equation}
\boxed{\;
b_{n} = a_{n}\left(e^{-\lambda_{n} D t_{c}} - 1\right) .
\;}
\label{eq:bn}
\end{equation}

No additional integrals are required: orthogonality of the eigenbasis transfers the
phase-one state directly. Since $\lambda_{n} > 0$, the inventory decays
exponentially after shutdown, with the late-time behaviour governed by
$\lambda_{1}$.

\section{Extension to the five-layer particle}

The full TRISO geometry comprises five concentric layers (kernel, buffer, IPyC, SiC,
OPyC) with radii
\begin{equation}
0 = r_{0} < r_{1} < r_{2} < r_{3} < r_{4} < r_{5} = R
\end{equation}
and piecewise constant diffusivities $D_{i}$. In each layer,
\begin{equation}
\frac{\partial c}{\partial t}
= \frac{D_{i}}{r^{2}}\,\frac{\partial}{\partial r}\!\left(r^{2}\frac{\partial c}{\partial r}\right)
+ S_{i} ,
\qquad i = 1,\dots,5 ,
\end{equation}
with the source confined to the kernel,
\begin{equation}
S_{1} = S_{0} ,
\qquad
S_{2} = S_{3} = S_{4} = S_{5} = 0 .
\end{equation}

\paragraph{Interface conditions.}
Both the concentration and the radial flux are continuous at each interface:
\begin{equation}
c\!\left(r_{i}^{-},t\right) = c\!\left(r_{i}^{+},t\right) ,
\qquad
D_{i}\,\frac{\partial c}{\partial r}\bigg|_{r = r_{i}^{-}}
= D_{i+1}\,\frac{\partial c}{\partial r}\bigg|_{r = r_{i}^{+}} ,
\qquad i = 1,\dots,4 .
\label{eq:interface}
\end{equation}

\paragraph{Composite spatial problem.}
Writing $D(r)$ for the piecewise diffusivity, the spatial problem becomes
\begin{equation}
\frac{d}{dr}\!\left(D(r)\, r^{2}\frac{d\phi}{dr}\right) + \lambda\, r^{2}\phi = 0 ,
\end{equation}
with time factor $e^{-\lambda t}$, the diffusivity now residing inside the
derivative rather than in the time equation. Flux continuity \eqref{eq:interface} is
exactly continuity of $D(r)\,r^{2}\,d\phi/dr$, which is the condition under which
self-adjointness, orthogonality
\begin{equation}
\int_{0}^{R} \phi_{n}\,\phi_{m}\, r^{2}\,dr = 0 ,
\qquad n \neq m ,
\end{equation}
and realness of the spectrum survive the material discontinuities.

\paragraph{Layer-wise solution.}
With $u = r\phi$ in each layer,
\begin{equation}
u^{(i)}(r) = \alpha_{i}\sin\!\left(\sqrt{\frac{\lambda}{D_{i}}}\;r\right)
+ \beta_{i}\cos\!\left(\sqrt{\frac{\lambda}{D_{i}}}\;r\right) ,
\qquad
\beta_{1} = 0
\end{equation}
by regularity at the origin. The remaining constants satisfy eight interface
equations from \eqref{eq:interface} together with the outer Robin condition, forming
a homogeneous linear system conveniently assembled as a product of $2 \times 2$
transfer matrices across the layers. Nontrivial solutions exist only where the
system determinant vanishes,
\begin{equation}
\det M(\lambda) = 0 ,
\end{equation}
which is the transcendental eigencondition generalizing \eqref{eq:eigcond}. The
steady-state split, the projection onto the eigenbasis, and the switch-off treatment
\eqref{eq:bn} then carry over unchanged from the single-layer case.

\end{document}
